\chapter{Representabilidad y D1}
\label{cap:d1}

Ya hay un $\Prov$ escrito en el lenguaje objeto. Falta comprobar que dice lo que pretende: que
cuando algo es demostrable de verdad, la teoría \emph{demuestra} que lo es.

\section{El enunciado}

\begin{teorema}[D1 — representabilidad positiva]
Si $\varphi$ tiene una demostración en el cálculo finitario, entonces la teoría demuestra
$\Prov(\codigo{\varphi})$.
\end{teorema}

\noindent En los dos cálculos, que es como el proyecto lo necesita:

\leanfrag{repr_pos}
\leanfrag{repr_pos_prf}

\selee{de que $\varphi$ tenga demostración finitaria se sigue que la teoría demuestra que
$\varphi$ tiene demostración.}

\section{Por qué D1 sale sin inducción}

Éste es el punto conceptual del capítulo, y conviene decirlo despacio porque es lo que separa D1 de
D3.

\begin{leccion}[Un testigo concreto es un cálculo]
La hipótesis de D1 no es «$\varphi$ es demostrable» en abstracto: es \textbf{una derivación
concreta}, un objeto finito que se tiene en la mano. La prueba consiste en \emph{traducirla} —
recorrer sus líneas, emitir el código de cada una, y comprobar que el verificador las acepta. Eso es
un cómputo finito sobre un objeto dado, y por eso \textbf{no necesita inducción de la teoría
objeto}: no hay que razonar sobre \emph{todas} las derivaciones, sólo sobre ésta.
\end{leccion}

Y está medido, no supuesto: el \emph{footprint} de D1 no cita ningún esquema de inducción — sólo el
ancla de codificación. Lo mismo vale para D2. En cambio Gödel~I y Gödel~II sí citan
\ident{ax_induction} y \ident{ax_list_induction}.

\begin{center}\small
\begin{tabularx}{\linewidth}{@{}l >{\raggedright\arraybackslash}X@{}}
\toprule
resultado & ¿usa inducción de la teoría objeto? \\
\midrule
\textbf{D1} \ident{repr_pos'_prf} & \textbf{no} — sólo el ancla de codificación \\
\textbf{D2} \ident{d2_prf} & \textbf{no} \\
Gödel I \ident{goedel_first_numeral} & sí \\
Gödel II \ident{goedel_second'} & sí (más el \ident{d3} pendiente) \\
\bottomrule
\end{tabularx}
\end{center}

\begin{muro}[Y la contrapartida, que es todo el resto del libro]
D1 es la $\Sigma_1$-completitud \textbf{externa}: la comprobamos nosotros, caso a caso. \textbf{D3}
es la misma propiedad \textbf{provable}: que la teoría demuestre, \emph{para toda} derivación, que
el verificador la acepta. Ahí sí hay que razonar sobre todas, y ahí sí hace falta inducción — y ése
es el muro que la Parte~IV cuenta.
\end{muro}

\section{El ancla de codificación}

Queda una pieza que el \emph{footprint} de D1 sí cita, y que conviene no esconder:
\ident{ax_axiomsCodeT_eq}. La regla del verificador que acepta un axioma de teoría necesita
comprobar que el código de la fórmula está en la lista de códigos de los axiomas — y esa lista,
escrita literalmente, es un término gigantesco. El axioma dice que el símbolo opaco
\ident{axiomsCodeT} \emph{es} esa lista, sin materializarla.

\begin{leccion}[Un ancla no es un postulado gödeliano]
Es una extensión \textbf{conservadora}: nombra un objeto que ya existe. No añade fuerza aritmética —
añade un nombre corto para un término largo, y evita que las pruebas tengan que manipularlo. Está
en el inventario de \doc{AXIOMS.md} como los axiomas 5 y 6, uno por cálculo.
\end{leccion}
